% GENERATED FILE. Edit canonical lessons, then run npm run generate:books.
% canonical-source-hash: fnv1a64:593ec452ec88364e
% canonical-lessons: ZH-C222-zhaoliang, ZH-C222-kesou, ZH-C222-sangzi, ZH-C222-dapenti, ZH-C222-nanshou

\chapter{Under the Weather}
\label{ch:zh-under-the-weather}
\hlchaptermodality{222}

\begin{chapteropening}
\textbf{By the end of this chapter:} \emph{I can say I caught a chill, I have a cough and a sore throat, I keep sneezing, and why I feel awful.}

Complete the last lesson of chapter 222: I can say I caught a chill, I have a cough and a sore throat, I keep sneezing, and why I feel awful.
\end{chapteropening}

\section[zháoliáng]{zháoliáng — to catch a chill}

\label{lesson:ZH-C222-zhaoliang}



\begin{quote}
Before the new one: say the Mandarin for below zero, then the Mandarin for gloves.
\end{quote}

\subsection*{You'll want to know: zháoliáng}
\emph{zháoliáng} — \textquotedblleft{}to catch a chill\textquotedblright{}.

\begin{quote}
\emph{zuótiān wǒ zháoliáng le} — I caught a chill yesterday.
\end{quote}

Its characters come later; for now it is a word for the ear.

\begin{cousinweb}
\emph{zháoliáng} is built from \emph{zháo} 'to be touched by' and \emph{liáng} 'cool', both inherited from Middle Chinese, \emph{drjak} and \emph{ljang} (Guangyun, in Baxter's transcription), which glosses the first 'to attach to' and the second 'thin' and 'cold'.
\end{cousinweb}

\begin{grammarlens}[title={What you've built}]
One more word for giving reasons.
\end{grammarlens}

\subsection*{Your turn}
\begin{itemize}
\raggedright
  \item \emph{Say it:} \emph{zháoliáng}
  \item \emph{Say it:} \emph{zháoliáng}, once more
  \item \emph{Say it:} \emph{zuótiān wǒ zháoliáng le} — I caught a chill yesterday.
  \item \emph{Recall:} say the Mandarin for below zero, then the Mandarin for gloves, then say \emph{zháoliáng} again
\end{itemize}

\begin{tcolorbox}[breakable,colback=teal!4,colframe=teal!35!black,title={Before you move on}]
What does zháoliáng mean? (\textquotedblleft{}To catch a chill\textquotedblright{}.) Say it once more.
\end{tcolorbox}
\section[késou]{késou — to cough}

\label{lesson:ZH-C222-kesou}



\begin{quote}
Before the new one: say the Mandarin for gloves, then the Mandarin for to catch a chill.
\end{quote}

\subsection*{You'll want to know: késou}
\emph{késou} — \textquotedblleft{}to cough\textquotedblright{}.

\begin{quote}
\emph{wǒ yǒudiǎnr késou} — I've got a bit of a cough.
\end{quote}

Its characters come later; for now it is a word for the ear.

\begin{cousinweb}
\emph{késou} is built from \emph{ké} and \emph{sou}. The second is inherited from Middle Chinese \emph{suwH} (Guangyun, in Baxter's transcription), listed with a character for coughing; the dictionary reads the first only as a small child's laugh, so how it came to mean coughing is not settled.
\end{cousinweb}

\begin{grammarlens}[title={What you've built}]
One more word for giving reasons.
\end{grammarlens}

\subsection*{Your turn}
\begin{itemize}
\raggedright
  \item \emph{Say it:} \emph{késou}
  \item \emph{Say it:} \emph{késou}, once more
  \item \emph{Say it:} \emph{wǒ yǒudiǎnr késou} — I've got a bit of a cough.
  \item \emph{Recall:} say the Mandarin for gloves, then the Mandarin for to catch a chill, then say \emph{késou} again
\end{itemize}

\begin{tcolorbox}[breakable,colback=teal!4,colframe=teal!35!black,title={Before you move on}]
What does késou mean? (\textquotedblleft{}To cough\textquotedblright{}.) Say it once more.
\end{tcolorbox}
\section[sǎngzi]{sǎngzi — throat; voice}

\label{lesson:ZH-C222-sangzi}



\begin{quote}
Before the new one: say the Mandarin for to catch a chill, then the Mandarin for to cough.
\end{quote}

\subsection*{You'll want to know: sǎngzi}
\emph{sǎngzi} — \textquotedblleft{}throat; voice\textquotedblright{}.

\begin{quote}
\emph{wǒ sǎngzi téng} — My throat hurts.
\end{quote}

Its characters come later; for now it is a word for the ear.

\begin{cousinweb}
\emph{sǎngzi} is built from \emph{sǎng} 'throat' and the noun ending \zh{子}, inherited from Middle Chinese \emph{tsiX} (Guangyun, in Baxter's transcription). No dictionary we use records an older form of the first, so its history is not settled.
\end{cousinweb}

\begin{grammarlens}[title={What you've built}]
One more word for giving reasons.
\end{grammarlens}

\subsection*{Your turn}
\begin{itemize}
\raggedright
  \item \emph{Say it:} \emph{sǎngzi}
  \item \emph{Say it:} \emph{sǎngzi}, once more
  \item \emph{Say it:} \emph{wǒ sǎngzi téng} — My throat hurts.
  \item \emph{Recall:} say the Mandarin for to catch a chill, then the Mandarin for to cough, then say \emph{sǎngzi} again
\end{itemize}

\begin{tcolorbox}[breakable,colback=teal!4,colframe=teal!35!black,title={Before you move on}]
What does sǎngzi mean? (\textquotedblleft{}Throat; voice\textquotedblright{}.) Say it once more.
\end{tcolorbox}
\section[dǎ pēntì]{dǎ pēntì — to sneeze}

\label{lesson:ZH-C222-dapenti}



\begin{quote}
Before the new one: say the Mandarin for to cough, then the Mandarin for throat.
\end{quote}

\subsection*{You'll want to know: dǎ pēntì}
\emph{dǎ pēntì} — \textquotedblleft{}to sneeze\textquotedblright{}.

\begin{quote}
\emph{gǎnmào de shíhou wǒ chángcháng dǎ pēntì} — When I've got a cold I sneeze a lot.
\end{quote}

Its characters come later; for now it is a word for the ear.

\begin{cousinweb}
\emph{dǎ pēntì} is built from \zh{打} 'to strike, to do' and \emph{pēntì} 'a sneeze', all inherited from Middle Chinese, \emph{taengX}, \emph{phwonH} and \emph{tejH} (Guangyun, in Baxter's transcription), which glosses the second 'to breathe out' and the third 'breath from the nose'. Neither dǎ nor the level tone of pēn is the regular outcome of the form listed.
\end{cousinweb}

\begin{grammarlens}[title={What you've built}]
One more word for giving reasons.
\end{grammarlens}

\subsection*{Your turn}
\begin{itemize}
\raggedright
  \item \emph{Say it:} \emph{dǎ pēntì}
  \item \emph{Say it:} \emph{dǎ pēntì}, once more
  \item \emph{Say it:} \emph{gǎnmào de shíhou wǒ chángcháng dǎ pēntì} — When I've got a cold I sneeze a lot.
  \item \emph{Recall:} say the Mandarin for to cough, then the Mandarin for throat, then say \emph{dǎ pēntì} again
\end{itemize}

\begin{tcolorbox}[breakable,colback=teal!4,colframe=teal!35!black,title={Before you move on}]
What does dǎ pēntì mean? (\textquotedblleft{}To sneeze\textquotedblright{}.) Say it once more.
\end{tcolorbox}
\section[nánshòu]{nánshòu — miserable; unwell}

\label{lesson:ZH-C222-nanshou}



\begin{quote}
Before the new one: say the Mandarin for throat, then the Mandarin for to sneeze.
\end{quote}

\subsection*{You'll want to know: nánshòu}
\emph{nánshòu} — \textquotedblleft{}miserable; unwell\textquotedblright{}.

\begin{quote}
\emph{yīnwèi gǎnmào le wǒ hěn nánshòu} — I feel awful because I've got a cold.
\end{quote}

Its characters come later; for now it is a word for the ear.

\begin{cousinweb}
\emph{nánshòu} is built from \emph{nán} 'hard' and \emph{shòu} 'to bear', both inherited from Middle Chinese, \emph{nan} and \emph{dzyuwX} (Guangyun, in Baxter's transcription), which glosses the first 'difficult' and the second 'to hold' and 'to receive'.
\end{cousinweb}

\begin{grammarlens}[title={What you've built}]
One more word for giving reasons.
\end{grammarlens}

\subsection*{Your turn}
\begin{itemize}
\raggedright
  \item \emph{Say it:} \emph{nánshòu}
  \item \emph{Say it:} \emph{nánshòu}, once more
  \item \emph{Say it:} \emph{yīnwèi gǎnmào le wǒ hěn nánshòu} — I feel awful because I've got a cold.
  \item \emph{Recall:} say the Mandarin for throat, then the Mandarin for to sneeze, then say \emph{nánshòu} again
\end{itemize}

\begin{tcolorbox}[breakable,colback=teal!4,colframe=teal!35!black,title={Before you move on}]
What does nánshòu mean? (\textquotedblleft{}Miserable; unwell\textquotedblright{}.) Say it once more.
\end{tcolorbox}
